\documentclass[12pt]{article}
\usepackage[utf8]{inputenc}
\usepackage[T1]{fontenc}
\usepackage{hyperref}
\title{Central Tunebo Language Resources for K-12 Curriculum Development}
\author{Compiled by Claude Sonnet 5 (claude-sonnet-5) for Andr\'as Kornai}
\date{}
\begin{document}
\maketitle

\section{Introduction}
Central Tunebo (ISO 639-3 tuf; Glottocode cent2150), also called Uw cuwa, U'wa, Cobaría Tunebo or Uwa-Tunebo, is a Chibchan language of Boyacá, Norte de Santander and Arauca in Colombia. Speaker counts in the searches vary, from about 1,800 to higher figures for the people group. Only a few resources were found.

\section{Primary Educational Resources}
No primers, textbooks or school programmes were found in the searches.

\section{Linguistic and Reference Materials}
\begin{itemize}
\item Bilingual Uw cuwa--Spanish dictionary on Webonary, with a shorter Spanish index: \url{https://www.webonary.org/tunebo/overview/introduction/?lang=en}
\end{itemize}

\section{Religious and Community Resources}
\begin{itemize}
\item Scripture Earth page for Tunebo, Central: online and audio Bible, New Testament PDF, some Old Testament selections, and Android Bible and songbook apps: \url{https://www.scriptureearth.org/00eng.php?iso=tuf}
\item Joshua Project profile of the Central Tunebo people (the Bible status given there is New Testament): \url{https://joshuaproject.net/people_groups/11279}
\end{itemize}

\section{Digital and Audio Resources}
The Scripture Earth page lists audio Scripture and mobile apps (see above). No other digital or audio resources were found.

\section{Curriculum Development Context}
The language is used by all in the community, which favors a mother-tongue curriculum. The searches showed no existing school programme. The Webonary dictionary and the Scripture materials could serve as the base for reading material and vocabulary. The community's own priorities, and any bilingual education work among the U'wa, were not covered by the searches and need further investigation.
\end{document}